\chapter{Réalisation de la plate-forme}
\label{ch:real}

\section{Introduction}

Les sources de la PaaS vivent sous~\textbf{\texttt{paas/frontend}}~: projet Node.js~(voir~\texttt{package.json}) où les~\emph{layouts}~\texttt{src/app/} composent les écrans du tableau de bord alors que~\texttt{src/app/api/}~(Route Handlers)~exposent notamment les~\texttt{POST~/~GET} pour les projets, les déploiements et les webhooks~(ex.~\texttt{/api/webhooks/github})~;, adaptateurs~\texttt{build-backend-*.ts}, accès registre/GitOps~(Harbor~\texttt{HARBOR\_*},~\texttt{GITOPS\_*},~\texttt{ARGOCD\_*}) ainsi que~\texttt{config/env.ts}.

Le développement local suit~\texttt{npm install ; npm run prisma:generate ; npm run dev}~(cf.~\texttt{paas/README.md})~; le mode~\texttt{DEVSECOPS\_ALLOW\_SIMULATION} permet de parcourir l'IU sans cluster.

Les pages qui suivent mentionnent encore un \emph{adaptateur Spring Boot}~(optionnel dans un autre module)~; elles ne prévalent pas sur le cœur Next.js défini aux chapitres~\ref{ch:archi}--\ref{ch:outils}.

\section{Objectifs de réalisation}

Le sprint associé au développement vise une IHM permettant l'onboarding des projets, le déclenchement ou le suivi des builds~(Jenkins ou Tekton)~et la vue d'ensemble des déploiements jusqu'à Argo CD.

\begin{figure}[h]
\centering
\includegraphics[width=\textwidth]{sprint4.png}
\caption{Vague de développement IHM ~/ intégrations}
\end{figure}

\subsection*{Composant serveur Java~(optionnel)}
Le dépôt peut inclure un service Spring Boot pour des points d'extension~(ex.\ proxys vers Jenkins ou Artifactory)~; le cœur utilisateur et les API métiers documentés dans la solution relèvent du \textbf{monorepo Next.js}. Lorsque des pages décrivent en détail un \texttt{JenkinsController}, elles illustrent ce mode d'intégration REST et non l'ensemble du produit.

\subsection*{Package Controller (exemple Jenkins, optionnel)}

Dans cette section, nous allons décrire la classe \textit{JenkinsController}, qui constitue un élément central de notre application backend.

\subsubsection*{Annotations :}

\begin{enumerate}
    \item \textbf{@RestController} : Cette annotation indique que la classe est un contrôleur REST, ce qui signifie qu’elle gère les requêtes HTTP et retourne directement des réponses sous forme de données (JSON généralement).

    \item \textbf{@RequestMapping("/jenkins")} : Cette annotation permet de définir le chemin de base pour tous les endpoints exposés dans ce contrôleur. Ainsi, toutes les API commencent par \textit{/jenkins}.

    \item \textbf{@CrossOrigin(origins = "*")} : Cette annotation autorise les requêtes HTTP provenant de différentes origines (CORS). Elle permet notamment au frontend Next.js de communiquer avec le service sans être bloqué par les restrictions de sécurité du navigateur.
\end{enumerate}

\subsubsection*{Rôle du JenkinsController :}

Le contrôleur joue le rôle d’intermédiaire entre l’interface utilisateur et les services backend. Il reçoit les requêtes HTTP envoyées par le frontend, les traite, puis appelle les services appropriés pour interagir avec le serveur Jenkins.

Par exemple :
\begin{itemize}
    \item Une requête \textbf{POST} permet de lancer un job Jenkins avec des paramètres spécifiques.
    \item Une requête \textbf{GET} permet de récupérer le statut d’un build Jenkins.
\end{itemize}

Le contrôleur délègue ensuite la logique métier au service, ce qui permet de respecter le principe de séparation des responsabilités et d’assurer une meilleure maintenabilité du code.
\subsubsection*{Méthodes du JenkinsController :}

\paragraph{1. Annotation @PostMapping("/jobwithparams/\{jobName\}") :}
Cette annotation indique que la méthode gère les requêtes HTTP POST avec le chemin \textit{/jobwithparams/\{jobName\}}. Le paramètre \textit{jobName} est un paramètre de chemin (PathVariable) extrait directement de l’URL.

\subsubsection*{Méthode : build}

\begin{enumerate}
    \item \textbf{Description :} Cette méthode permet d’exécuter un job Jenkins en envoyant une requête HTTP POST avec des paramètres spécifiques.

    \item \textbf{Chemin de l’API :} \textit{/jenkins/jobwithparams/\{jobName\}}

    \item \textbf{Méthode HTTP :} POST

    \item \textbf{Paramètres :}
    \begin{itemize}
        \item \textbf{jobName (PathVariable)} : Nom du job Jenkins à exécuter.
        \item \textbf{params (RequestParam)} : Paramètres nécessaires à l’exécution du job.
    \end{itemize}

    \item \textbf{Fonctionnalité :}  
    Cette méthode appelle le service Jenkins pour déclencher l’exécution du job avec les paramètres fournis.
\end{enumerate}

\subsubsection*{Méthode : getJobStatus}

\begin{enumerate}
    \item \textbf{Description :} Cette méthode permet de récupérer le statut d’un job Jenkins via une requête HTTP GET.

    \item \textbf{Chemin de l’API :} \textit{/jenkins/\{jobName\}/getJobstatus}

    \item \textbf{Méthode HTTP :} GET

    \item \textbf{Paramètres :}
    \begin{itemize}
        \item \textbf{jobName (PathVariable)} : Nom du job Jenkins dont on souhaite obtenir le statut.
    \end{itemize}

    \item \textbf{Fonctionnalité :}  
    Cette méthode interroge le service Jenkins pour récupérer le statut du job spécifié, puis retourne une réponse HTTP contenant ce statut.
\end{enumerate}

\subsection*{Package Service :}

Dans cette section, nous décrivons la classe \textit{JenkinsService}, responsable de la logique métier et de la communication avec le serveur Jenkins.

\begin{enumerate}
    \item \textbf{@Service :}  
    Cette annotation indique que la classe est un composant de service dans Spring, utilisé pour encapsuler la logique métier.

    \item \textbf{Constructeur JenkinsService :}  
    Le constructeur initialise une instance de \textit{JenkinsServer} en utilisant :
    \begin{itemize}
        \item L’URL du serveur Jenkins
        \item Le nom d’utilisateur
        \item Le token API
    \end{itemize}
    Ces valeurs sont injectées via les annotations \textit{@Value} et permettent d’établir une connexion sécurisée avec le serveur Jenkins.
\end{enumerate}
\subsubsection*{Implémentation du JenkinsService :}

\paragraph{Méthode : runJobWithParams}

Cette méthode permet d’exécuter un job Jenkins avec des paramètres spécifiques et de gérer son déclenchement côté serveur.

\begin{itemize}
    \item \textbf{Entrées :}
    \begin{itemize}
        \item Nom du job Jenkins
        \item Map de paramètres (clé/valeur)
    \end{itemize}

    \item \textbf{Étapes :}
    \begin{itemize}
        \item Récupération de l’objet \textit{Job} correspondant depuis le serveur Jenkins.
        \item Déclenchement d’un nouveau build avec les paramètres fournis.
    \end{itemize}
\end{itemize}

\paragraph{Méthode : getJobStatus}

Cette méthode permet de récupérer le statut d’un job Jenkins donné.

\begin{itemize}
    \item \textbf{Entrée :}
    \begin{itemize}
        \item Nom du job Jenkins
    \end{itemize}

    \item \textbf{Étapes :}
    \begin{itemize}
        \item Récupération de l’objet \textit{Job} depuis Jenkins.
        \item Récupération du dernier build associé à ce job.
        \item Extraction du résultat du build (SUCCESS, FAILURE, etc.).
        \item Retour du statut sous forme de chaîne de caractères.
    \end{itemize}
\end{itemize}

\subsection*{Configuration application.properties :}

Dans le fichier \textit{application.properties}, nous définissons les paramètres nécessaires à la connexion avec le serveur Jenkins.

\begin{itemize}
    \item \textbf{jenkins.url} :  
    URL du serveur Jenkins auquel l’application Spring Boot se connecte.

    \item \textbf{jenkins.username} :  
    Nom d’utilisateur utilisé pour l’authentification avec Jenkins.

    \item \textbf{jenkins.api.token} :  
    Token API utilisé pour sécuriser l’authentification.  
    Grâce à ces paramètres, l’application peut communiquer de manière sécurisée avec Jenkins.

\end{itemize}

\paragraph{Remarque :}  
Une fois configurée, l’application est accessible via une URL de type :  
\textit{http://localhost:8080/myapp}
\subsection*{Configuration SSL :}

Dans le cas où l’API Jenkins est sécurisée via HTTPS, il est nécessaire d’importer le certificat dans le keystore de la JVM afin d’établir une connexion sécurisée.

Pour cela, nous utilisons la commande suivante :

\begin{verbatim}
keytool -import -alias mycert \
-keystore /chemin/vers/mon/magasin-de-cles.jks \
-file /chemin/vers/le/certificat.crt \
-storepass mot_de_passe -keypass mot_de_passe
\end{verbatim}

Cette étape permet à l’application Spring Boot de faire confiance au certificat du serveur Jenkins et d’éviter les erreurs SSL lors des appels HTTPS.

\subsection*{Dépendance Jenkins :}

Pour interagir avec Jenkins depuis notre application, nous utilisons la dépendance suivante :

\begin{verbatim}
<dependency>
    <groupId>org.jenkins-ci.main</groupId>
    <artifactId>jenkins-core</artifactId>
    <version>2.289.3</version>
</dependency>
\end{verbatim}

Cette dépendance permet d’accéder aux fonctionnalités principales de Jenkins ainsi qu’aux classes nécessaires pour manipuler les jobs, builds et autres éléments du serveur Jenkins.
\subsection*{6.2.4 Développement du Frontend :}

\subsubsection*{Description du fonctionnement :}

Le frontend de l’application est développé en utilisant le framework \textbf{Angular}. Il fournit une interface utilisateur intuitive permettant d’interagir facilement avec les fonctionnalités de gestion des jobs Jenkins.

Le frontend est composé principalement des éléments suivants :

\subsubsection*{Formulaire de lancement du build}

Le formulaire permet à l’utilisateur de renseigner les paramètres nécessaires pour lancer un job Jenkins. Il contient plusieurs champs de saisie correspondant aux paramètres attendus (nom du projet, tag, type de build, etc.).

Lorsque l’utilisateur clique sur le bouton \textbf{"Build"} :
\begin{itemize}
    \item Les données saisies sont collectées depuis le formulaire.
    \item Une requête HTTP \textbf{POST} est envoyée vers le backend.
    \item Le backend déclenche l’exécution du job Jenkins avec les paramètres fournis.
\end{itemize}

\paragraph{Fonctionnement du code :}
\begin{itemize}
    \item Vérification de la validité du formulaire.
    \item Création d’un objet contenant les données saisies.
    \item Ajout de ces données dans un tableau local.
    \item Sauvegarde dans le \textit{localStorage}.
    \item Appel de la fonction qui déclenche le build Jenkins.
\end{itemize}

\subsubsection*{Tableau de données}

Le frontend affiche également un tableau récapitulatif contenant les informations des builds exécutés.

\begin{itemize}
    \item Les données du formulaire sont affichées après soumission.
    \item Chaque ligne représente un job Jenkins lancé.
    \item Les paramètres utilisés sont conservés pour consultation ultérieure.
\end{itemize}

Ce tableau permet d’avoir une vue claire et organisée des builds Jenkins exécutés précédemment, facilitant ainsi le suivi et la gestion des opérations.
\subsubsection*{Champ de récupération du statut du build}

Après soumission du formulaire via le bouton \textbf{"Submit"}, une requête HTTP \textbf{POST} est envoyée au backend afin de déclencher le build Jenkins avec les paramètres fournis.

\begin{itemize}
    \item Le frontend affiche un \textbf{spinner} pour indiquer que le build est en cours d’exécution.
    \item Une fois le build lancé, une requête HTTP \textbf{GET} est envoyée automatiquement pour récupérer le statut du build.
    \item Le backend interagit avec Jenkins pour obtenir ce statut.
\end{itemize}

Lorsque le statut est récupéré :
\begin{itemize}
    \item Il est automatiquement mis à jour dans la colonne \textbf{"Status"} du tableau.
    \item Si le statut est \textbf{"pending"}, le spinner reste affiché.
    \item Une fois le build terminé, le statut final (SUCCESS, FAILURE, etc.) est affiché.
\end{itemize}

Cette approche permet :
\begin{itemize}
    \item Une mise à jour en temps réel du statut des builds.
    \item Une meilleure expérience utilisateur.
    \item Une visibilité claire sur l’état des jobs Jenkins.
\end{itemize}

\subsubsection*{BuildService}

Le service \textbf{BuildService} est responsable de la communication entre le frontend et le backend.

\begin{itemize}
    \item Il utilise le service \textbf{HttpClient} d’Angular pour effectuer les appels HTTP.
\end{itemize}

\paragraph{Méthodes principales :}

\begin{itemize}
    \item \textbf{buildJob} :  
    Envoie une requête HTTP \textbf{POST} vers l’API backend pour déclencher un job Jenkins avec les paramètres fournis.  
    Les paramètres sont transmis sous forme de \textbf{HttpParams}.

    \item \textbf{getJobStatus} :  
    Envoie une requête HTTP \textbf{GET} pour récupérer le statut d’un job Jenkins spécifique.
\end{itemize}

\paragraph{Configuration :}

\begin{itemize}
    \item \textbf{API\_URL} :  
    Contient l’URL de base de l’API backend, généralement définie dans les fichiers d’environnement Angular.
\end{itemize}
\subsection*{Résultat}

Dans cette section, nous présentons une application fonctionnelle répondant aux besoins de gestion des travaux Jenkins. L’interface homme-machine (IHM) développée permet aux utilisateurs de lancer des jobs Jenkins avec des paramètres personnalisés et de suivre en temps réel le statut des builds.

L’application offre les fonctionnalités suivantes :
\begin{itemize}
    \item Lancement de jobs Jenkins avec des paramètres dynamiques.
    \item Suivi en temps réel de l’état des builds (pending, success, failure).
    \item Affichage des résultats dans un tableau structuré pour une meilleure lisibilité.
    \item Interface intuitive facilitant l’interaction utilisateur.
\end{itemize}

\begin{figure}[H]
    \centering
    \includegraphics[width=0.9\textwidth]{images/ihm-app.png}
    \caption{Interface utilisateur de l'application}
\end{figure}

\subsection*{Conteneurisation des deux applications}

La conteneurisation représente un élément essentiel de notre projet DevSecOps. Elle permet de standardiser le déploiement et d’assurer une portabilité ainsi qu’une reproductibilité des environnements.

Nous avons adopté Docker pour conteneuriser les deux applications principales :
\begin{itemize}
    \item \textbf{Application Backend} développée avec Spring Boot
    \item \textbf{Application Frontend} développée avec Angular
\end{itemize}

Cette approche offre plusieurs avantages :
\begin{itemize}
    \item Isolation des environnements d’exécution
    \item Déploiement simplifié
    \item Scalabilité et maintenance facilitées
\end{itemize}

Dans les sections suivantes, nous détaillons les fichiers \textbf{Dockerfile} utilisés pour chaque application.
\subsection*{Dockerfile pour l'application Backend}

Le conteneur backend est basé sur une image OpenJDK provenant d’un registre privé. Les étapes de construction sont les suivantes :

\begin{itemize}
    \item \textbf{Étape 1 :} Utilisation d'une image de base OpenJDK.
    \item \textbf{Étape 2 :} Définition d'une variable d'argument \texttt{JAR\_NAME}.
    \item \textbf{Étape 3 :} Ajout du fichier JAR généré dans le conteneur.
    \item \textbf{Étape 4 :} Exposition du port d'écoute.
    \item \textbf{Étape 5 :} Définition de la commande de démarrage.
\end{itemize}

\begin{lstlisting}[language=Dockerfile, caption=Dockerfile du Backend]
FROM harbor-stage-abylsen.com/abylsen/openjdk:17-bullseye

ARG JAR_NAME

ADD target/${JAR_NAME} app.jar

EXPOSE 8080

ENTRYPOINT ["java", "-jar", "app.jar"]
\end{lstlisting}

\subsection*{Dockerfile pour l'application Frontend}

Le conteneur frontend repose sur une approche multi-stage permettant d’optimiser la taille de l’image finale.

\begin{itemize}
    \item \textbf{Étape 1 :} Utilisation d'une image Node.js depuis le registre privé.
    \item \textbf{Étape 2 :} Définition du répertoire de travail.
    \item \textbf{Étape 3 :} Copie des fichiers \texttt{package.json} et \texttt{package-lock.json}.
    \item \textbf{Étape 4 :} Installation des dépendances.
    \item \textbf{Étape 5 :} Copie du code source.
    \item \textbf{Étape 6 :} Build de l’application Angular.
    \item \textbf{Étape 7 :} Utilisation d'une image Nginx légère.
    \item \textbf{Étape 8 :} Copie des fichiers buildés vers Nginx.
    \item \textbf{Étape 9 :} Ajout de la configuration Nginx.
    \item \textbf{Étape 10 :} Exposition du port.
    \item \textbf{Étape 11 :} Commande de démarrage.
\end{itemize}

\begin{lstlisting}[language=Dockerfile, caption=Dockerfile du Frontend]
 Stage 1: Build Angular app
FROM harbor-stage-abylsen.com/abylsen/node:18 AS build

WORKDIR /app

COPY package*.json ./

RUN npm install

COPY . .

RUN npm run build --prod

 Stage 2: Serve with Nginx
FROM nginx:alpine

COPY --from=build /app/dist /usr/share/nginx/html

COPY nginx.conf /etc/nginx/nginx.conf

EXPOSE 80

CMD ["nginx", "-g", "daemon off;"]
\end{lstlisting}
\subsection*{Dockerfile du Frontend}

Le Dockerfile du frontend suit une approche multi-stage afin d’optimiser la taille de l’image finale et de séparer la phase de build de la phase de déploiement.

\begin{lstlisting}[language=Dockerfile, caption=Dockerfile du Frontend]
# Stage 1: Build Angular app
FROM harbor-stage-abylsen.com/abylsen/node:16 AS build

WORKDIR /app

COPY package*.json ./
RUN npm install

COPY . .
RUN npm run build

# Stage 2: Serve with Nginx
FROM nginx:1.21-alpine

COPY --from=build /app/dist /usr/share/nginx/html
COPY nginx.conf /etc/nginx/nginx.conf

EXPOSE 80

CMD ["nginx", "-g", "daemon off;"]
\end{lstlisting}

\subsection*{Chart Helm des deux applications}

Une chart Helm est un ensemble de fichiers YAML permettant de définir, paramétrer et déployer une application sur un cluster Kubernetes. Elle facilite la gestion du cycle de vie des applications (déploiement, mise à jour, rollback).

Une chart Helm est généralement structurée comme suit :

\begin{itemize}
    \item \textbf{Chart.yaml} : Contient les métadonnées de la chart (nom, version, description).
    \item \textbf{README.md} : Fournit des informations détaillées sur l’utilisation de la chart.
    \item \textbf{values.yaml} : Définit les valeurs par défaut utilisées dans les templates.
    \item \textbf{templates/} : Contient les fichiers YAML des ressources Kubernetes.
\end{itemize}

Dans notre cas, les principaux fichiers utilisés sont :

\begin{itemize}
    \item \textbf{deployment.yaml} : Définit la configuration du déploiement de l’application (pods, replicas, images, etc.).
    \item \textbf{service.yaml} : Définit un service Kubernetes pour exposer l’application.
    \item \textbf{\_helpers.tpl} : Contient des fonctions et variables réutilisables dans les templates.
\end{itemize}

Cette organisation permet de standardiser et d’automatiser le déploiement des applications backend et frontend dans notre environnement Kubernetes.
\subsection*{Revue du Sprint}

Le sprint s’est achevé avec succès après la réalisation de l’ensemble des tâches planifiées. L’équipe Scrum s’est réunie afin de valider le bon fonctionnement de l’application développée, en s’assurant que toutes les fonctionnalités implémentées répondent aux exigences définies.

Cette revue a également permis :
\begin{itemize}
    \item de vérifier la qualité des livrables produits,
    \item d’identifier les éventuels axes d’amélioration,
    \item de préparer efficacement les objectifs du prochain sprint.
\end{itemize}

\begin{figure}[H]
    \centering
    \includegraphics[width=0.9\textwidth]{images/sprint4.png}
    \caption{Fin du Sprint 4}
\end{figure}

\subsection*{Conclusion}

Dans ce chapitre, nous avons réussi à concevoir et développer une application complète permettant aux utilisateurs de lancer et de suivre des travaux Jenkins via une interface intuitive.

Les principaux apports de ce sprint incluent :
\begin{itemize}
    \item La mise en place d’une interface utilisateur facilitant l’interaction avec Jenkins,
    \item L’intégration d’un backend robuste permettant de piloter les jobs Jenkins,
    \item L’automatisation du suivi des builds en temps réel.
\end{itemize}

Ces réalisations constituent une base solide pour la prochaine étape du projet, qui portera sur la mise en place et l’automatisation complète du pipeline Jenkins.
